
\documentclass[10pt,a4paper,english]{report}
\usepackage{amssymb}




%%%%%%%%%%%%%%%%%%%%%%%%%%%%%%%%%%%%%%%%%%%%%%%%%%%%%%%%%%%%%%%%%%%%%%%%%%%%%%%%%%%%%%%%%%%%%%%%%%%%
%\usepackage{makeidx}
%\usepackage[T1]{fontenc}
%%\usepackage[latin9]{inputenc}
\usepackage{amsmath}
\usepackage{graphicx}
\usepackage{graphics}
\usepackage[dvips]{epsfig}
\usepackage{esint}
\usepackage[authoryear]{natbib}
\usepackage{rotating}
\usepackage{relsize}
\usepackage{Sweave}
\usepackage{xcolor}
\usepackage{babel}
\usepackage{babel}
\usepackage{hyperref}



\begin{document}
\title{Computer Exercise 2: Generalized Linear Models}
\author{Lars Ronnegard}
\maketitle
\section*{Introduction}
The objective of this exercise is to explore generalized linear models using some Poisson data and some binary data. Thereafter you will explore how to estimate the residual variance from a linear model using a Gamma distribution.


\section*{Exercise 1}
There are 15 observations from a Poisson distribution
\begin{Schunk}
\begin{Sinput}
y <-  c(4, 5, 4, 5, 3, 1, 2, 3, 4, 8, 0, 1, 4, 0, 2)
\end{Sinput}
\end{Schunk}

These data were generated from $Po(\lambda)$ and your task is to guess the value of $\lambda$. 

\section*{Exercise 2}
Use the glm() function in R to estimate $\lambda$. Do you get the same estimate as your guess in Exercise 1? \newline
Plot diagnostics by plotting the output from the glm() function and try to understand what the plots show.
Your code should look something like:
\begin{Schunk}
\begin{Sinput}
mod1 <- glm(....)
plot(mod1)
\end{Sinput}
\end{Schunk}

\section*{Exercise 3}
In a case-control study (with binary response variable $y$)  10 healthy and 10 sick individuals have been selected for a study. They were genotyped for a SNP suspected to be close to the causative mutation. The recorded genotypes were coded as 0, 1 or 2 (with 1=heterozygote) and saved as a vector $x$.    
\begin{Schunk}
\begin{Sinput}
y <- c(0,0,0,0,0,0,0,0,0,0,1,1,1,1,1,1,1,1,1,1)
x <- c(0,1,1,0,0,1,0,1,0,1,2,2,0,1,1,2,1,1,2,1)
\end{Sinput}
\end{Schunk}
Use $y$ as response in a binomial GLM to check whether the SNP-effect is significant.

\section*{Exercise 4}
Let us have a look at the savings data again.
\begin{Schunk}
\begin{Sinput}
library(faraway)
data(savings) #See Faraway page 29
?savings
head(savings)
summary(savings)
g <- lm(sr ~ pop15 + pop75 + dpi + ddpi, data=savings)
summary(g)
\end{Sinput}
\end{Schunk}
What is the residual variance? \newline
Now we will compute the residual variance using a Gamma distribution. First the residuals and hat values are extracted:
\begin{Schunk}
\begin{Sinput}
res1 <- residuals(g)
hv1 <- hatvalues(g)
\end{Sinput}
\end{Schunk}
Check that the following code gives an estimate of the residual variance.
\begin{Schunk}
\begin{Sinput}
response1 <- (res1^2)/(1-hv1)
weight1 <- (1-hv1)/2
glm1 <- glm(response1 ~ 1, weights = weight1, family=Gamma(link=log))
summary(glm1)
\end{Sinput}
\end{Schunk}
Compute a 95 \% confidence interval for the residual variance.

\section*{Exercise 4 (if you get time)}
Implement your own function in R for fitting a Poisson GLM using an Iterative Rewighted Least Squares algorithm.

\end{document}


